% Options for packages loaded elsewhere
\PassOptionsToPackage{unicode}{hyperref}
\PassOptionsToPackage{hyphens}{url}
\documentclass[
  a4paper,
]{article}
\usepackage{xcolor}
\usepackage[margin=25mm]{geometry}
\usepackage{amsmath,amssymb}
\setcounter{secnumdepth}{-\maxdimen} % remove section numbering
\usepackage{iftex}
\ifPDFTeX
  \usepackage[T1]{fontenc}
  \usepackage[utf8]{inputenc}
  \usepackage{textcomp} % provide euro and other symbols
\else % if luatex or xetex
  \usepackage{unicode-math} % this also loads fontspec
  \defaultfontfeatures{Scale=MatchLowercase}
  \defaultfontfeatures[\rmfamily]{Ligatures=TeX,Scale=1}
\fi
\usepackage{lmodern}
\ifPDFTeX\else
  % xetex/luatex font selection
\fi
% Use upquote if available, for straight quotes in verbatim environments
\IfFileExists{upquote.sty}{\usepackage{upquote}}{}
\IfFileExists{microtype.sty}{% use microtype if available
  \usepackage[]{microtype}
  \UseMicrotypeSet[protrusion]{basicmath} % disable protrusion for tt fonts
}{}
\makeatletter
\@ifundefined{KOMAClassName}{% if non-KOMA class
  \IfFileExists{parskip.sty}{%
    \usepackage{parskip}
  }{% else
    \setlength{\parindent}{0pt}
    \setlength{\parskip}{6pt plus 2pt minus 1pt}}
}{% if KOMA class
  \KOMAoptions{parskip=half}}
\makeatother
\usepackage{longtable,booktabs,array}
\newcounter{none} % for unnumbered tables
\usepackage{calc} % for calculating minipage widths
% Correct order of tables after \paragraph or \subparagraph
\usepackage{etoolbox}
\makeatletter
\patchcmd\longtable{\par}{\if@noskipsec\mbox{}\fi\par}{}{}
\makeatother
% Allow footnotes in longtable head/foot
\IfFileExists{footnotehyper.sty}{\usepackage{footnotehyper}}{\usepackage{footnote}}
\makesavenoteenv{longtable}
\setlength{\emergencystretch}{3em} % prevent overfull lines
\providecommand{\tightlist}{%
  \setlength{\itemsep}{0pt}\setlength{\parskip}{0pt}}
\usepackage{bookmark}
\IfFileExists{xurl.sty}{\usepackage{xurl}}{} % add URL line breaks if available
\urlstyle{same}
\hypersetup{
  hidelinks,
  pdfcreator={LaTeX via pandoc}}

\author{}
\date{}

\begin{document}

\section{\texorpdfstring{Symmetry Generation from Zero Closure, Finite
Order, and Self-Consistent Geometry --- The Single External Parameter
\(N\) and the Remaining Tasks of Generalization and
Dynamics}{Symmetry Generation from Zero Closure, Finite Order, and Self-Consistent Geometry --- The Single External Parameter N and the Remaining Tasks of Generalization and Dynamics}}\label{symmetry-generation-from-zero-closure-finite-order-and-self-consistent-geometry-the-single-external-parameter-n-and-the-remaining-tasks-of-generalization-and-dynamics}

\textbf{Author:} Noriaki Kihara**ORCID:** 0009-0004-6753-4020**Date:**
August 20, 2026**Version DOI:** 10.5281/zenodo.22028073**Concept DOI:**
10.5281/zenodo.22028072**Series:** ``Generative Structure of
Dimensions'' series --- Symmetry Derivation from Closure Axioms, public
version v1.0**License:** CC BY 4.0

\begin{quote}
\textbf{Subject}: This paper organizes how far geometry, symmetry,
statistics, and readout structure can be derived from an axiomatic
framework that, apart from the parts not yet derived or generalized,
contains no explicit continuous tuning parameters and uses only a single
discrete integer \(N\) as its external parameter. While the five complex
degrees of freedom are derived rigorously, the self-consistent selection
rule among readout sectors including the \(3+2\) decomposition, the
generalization of some structures, and the derivation of complete
dynamics remain the principal tasks ahead.
\end{quote}

\subsection{Citation Policy and the Classification of
Self-Hypotheses}\label{citation-policy-and-the-classification-of-self-hypotheses}

This paper verifies correspondence with known mathematics and known
physics primarily through external literature. However, among the
constructions used here, those already published as the author's own
numerical experiments and constructive derivations, and directly
required by the derivation chain of this paper, are adopted as
self-references --- the following three papers only.

\begin{enumerate}
\def\labelenumi{\arabic{enumi}.}
\tightlist
\item
  Noriaki Kihara, \textbf{``The Periodic Table of Waves v2 --- Particle
  Classification by Winding-Number Address and Observation Clock, and
  the Unification of Mass, Lifetime, and Splitting via the Clock Field
  \(\omega(x)\)''}, Zenodo, Concept DOI: 10.5281/zenodo.21830706.
\item
  Noriaki Kihara, \textbf{``Two-Layer Separation of Waves and Fields ---
  Unification of Gauge Fields and Gravitational Fields via a Universal
  Field-Readout Function''}, Zenodo, DOI: 10.5281/zenodo.21832257.
\item
  Noriaki Kihara, \textbf{``Zero Closure Was Four-Dimensional ---
  `Central Projection' Survives Even in the Complex World''}, Zenodo,
  Version DOI: 10.5281/zenodo.21902806, Concept DOI:
  10.5281/zenodo.21902805.
\end{enumerate}

Results depending on these three papers are never conflated with
theorems established by external theory; they are explicitly marked as

\[
\boxed{\text{previously derived results by self-hypothesis and self-construction}}
\]

throughout.

In particular, what this paper imports from ``The Periodic Table of
Waves v2'' is the \(Z_2\) decomposition into odd and even harmonics, the
sign under half-period translation, the reflectance determined by the
odd-harmonic ratio, the clock covering degree, the antipodal two-point
structure, and the Bose/Fermi/mixed statistical classification. These
are not placed in the ``underived'' category here; they are classified
as \textbf{derived, constructed, and numerically confirmed as
self-hypotheses}.

Likewise, the separation of the wave-existence layer from the
observational readout layer, and the placement of clock, curvature, and
gravity on the readout side, imported from ``Two-Layer Separation of
Waves and Fields'', are treated as \textbf{previously derived results as
self-hypotheses}.

By this citation policy, the paper clearly separates ``known results
from external literature'', ``rigorous derivations within this paper'',
``previously derived self-hypothesis results originating in the three
papers above'', and ``unresolved generalization tasks''.

\subsection{Reconstructing Curvature, Spacetime Signature, Quantization,
and Internal Symmetry from a Zero Closure Containing Unobservable
Complex
Axes}\label{reconstructing-curvature-spacetime-signature-quantization-and-internal-symmetry-from-a-zero-closure-containing-unobservable-complex-axes}

\textbf{Version:} public v1.0\\
\textbf{Date:} 2026-08-20\\
\textbf{Citation policy:} Self-citation is kept to the minimum. Among
the author's published papers, only the three primary sources whose
results are treated here as ``already derived, constructed, and
numerically confirmed'' are self-cited; external literature is used for
all other mathematical and mathematical-physics background, known
theorems, and comparison with prior work.

\begin{center}\rule{0.5\linewidth}{0.5pt}\end{center}

\subsection{Abstract}\label{abstract}

The central feature of this paper is that, within the range of what has
been derived, no continuous free parameters tuned from outside are
introduced, and the only externally specified quantity is a discrete
integer \(N\). Extending some of the already constructed and numerically
confirmed results into general theorems, and deriving complete dynamics
including local interactions in closed form from the axiom system, are
explicitly left as the principal tasks beyond this paper.

This paper examines whether the many symmetries and geometric structures
introduced separately in modern theoretical physics can be reconstructed
in a unified way from a small number of closure principles.

The central axiom is, throughout,

\[
\boxed{
\mathrm{A1}:\quad
\sum_{n=1}^{M} x_n^2=0,
\qquad x_n\in\mathbb C
}
\]

In this paper, the curvature radius \(R\), the time axis \(t\), and the
internal axis \(Q\) are not added outside this equation. They are all
read as components of the complex \(x_n\) that carry a sign different
from the observable real directions.

For example, for the four components

\[
(x,y,z,iR)
\]

A1 gives

\[
x^2+y^2+z^2+(iR)^2=0
\]

that is,

\[
\boxed{
x^2+y^2+z^2=R^2
}
\]

Therefore \(\sum x_n^2=R^2\) is neither an axiom separate from A1 nor a
different level set. It is \textbf{the real-form display of A1 with one
component read as the unobservable imaginary direction \(iR\)}.

Similarly, for

\[
(x,y,z,it,iR,iQ)
\]

we get

\[
x^2+y^2+z^2-t^2-R^2-Q^2=0.
\]

Thus the indefinite signature of spacetime, the curvature radius, and
the internal axis all emerge from the same quadratic form as real-form
displays of complex axes, without breaking the zero closure.

This is the central point of the paper.\\
Curvature need not be added as a separate field. The real-form display
of an unobservable complex axis appears on the observable side as a
curvature radius. By the same principle, the negative sign of the time
axis is not introduced from an external Lorentz metric: reading the
unobservable axis \(t\) as \(it\) produces \((it)^2=-t^2\). Hence \(R\),
\(t\), and \(Q\) are unobservable complex axes of the same kind, and the
origin of the negative signs is one and the same.

Furthermore,

\[
\boxed{
\mathrm{A2}:\quad U^N=I
}
\]

provides finite cyclicity, discrete phases, cyclotomic eigenvalues, and
Born-type squared weights;

\[
\boxed{
\mathrm{A3}:\quad
\text{all pairwise distances of }N\text{ vertices are consistent as a simplex}
}
\]

converts the relational set into geometry; and

\[
\boxed{
\mathrm{A4}:\quad X=\mathcal F(X)
}
\]

provides self-consistent fixed points, stabilizers, and symmetry
selection.

From these few principles, this paper classifies in the forward
direction the orthogonal groups, cyclic groups, permutation groups,
symplectic structure, Lorentz-type signature, constant-curvature
geometry, conformal structure, simplicial chain complexes, fixed-point
symmetries, and further the automorphism group of the
five-degree-of-freedom structure obtained from this axiom system, and
examines whether the internal symmetry groups used in the Standard Model
appear as a result.

\begin{center}\rule{0.5\linewidth}{0.5pt}\end{center}

\section{1. The Starting Point Is a Single Zero
Closure}\label{the-starting-point-is-a-single-zero-closure}

The first axiom of this paper is

\[
\boxed{
\sum_n x_n^2=0
}
\]

and nothing else.

What matters is that

\[
x_n\in\mathbb C
\]

When this equation is written as

\[
\sum_{\rm visible} x_a^2
=
R^2
\]

the right-hand side \(R^2\) has not been introduced as a new conserved
quantity.

The original equation contains a component

\[
x_R=iR
\]

and since

\[
x_R^2=(iR)^2=-R^2
\]

we have merely transposed

\[
\sum_{\rm visible}x_a^2-R^2=0
\]

Therefore

\[
\boxed{
\sum_{\rm visible}x_a^2=R^2
}
\]

is the decomposition display of

\[
\boxed{
\sum_nx_n^2=0
}
\]

into observable and unobservable components.

\begin{center}\rule{0.5\linewidth}{0.5pt}\end{center}

\subsection{Anonymity Does Not Mean Arbitrariness --- Selection of
Symmetry Sectors by Strong
Constraints}\label{anonymity-does-not-mean-arbitrariness-selection-of-symmetry-sectors-by-strong-constraints}

In this construction, no physical labels such as ``space'', ``time'',
``curvature'', ``charge'', or ``color'' are assigned in advance to the
fundamental relational components. The same zero closure can therefore
be read as multiple contractions, refinements, and decompositions
depending on the observation map.

For example, collecting

\[
r^2=x^2+y^2+z^2
\]

we may read

\[
r^2=t^2+R^2+Q^2
\]

Contracting further as

\[
R'^2=t^2+R^2+Q^2
\]

gives the three-dimensional ellipsoidal/spherical readout

\[
r^2=R'^2,
\qquad
x^2+y^2+z^2=R'^2
\]

On the other hand, the Lorentz-type decomposition

\[
x^2+y^2+z^2-t^2=R^2+Q^2
\]

is also possible, and refining the internal readout as

\[
Q^2=Q_1^2+Q_2^2+Q_3^2
\]

we may read

\[
x^2+y^2+z^2-t^2
=
R^2+Q_1^2+Q_2^2+Q_3^2
\]

However,

\[
\boxed{\text{anonymity}\neq\text{arbitrariness}}
\]

The existence of multiple readout candidates does not mean that physical
interpretations can be freely invented to fit any observation. The
states, decompositions, and readouts permitted in this construction must
simultaneously satisfy the strong constraints imposed independently in
this paper: zero closure, finite-order recurrence, simplex consistency,
harmonic structure, and self-consistent fixed points.

Conceptually, the allowed set is expressed as

\[
\boxed{
\mathcal S_{\rm allowed}
=
\mathcal S_{\rm closure}
\cap
\mathcal S_{\rm recurrence}
\cap
\mathcal S_{\rm simplex}
\cap
\mathcal S_{\rm harmonic}
\cap
\mathcal S_{\rm self-consistent}
}
\]

Therefore, even if the space of input candidates is large, the allowed
solution set that actually remains is restricted by this intersection.

The characteristic of this model is not the production of many
structures through free interpretation. On the contrary, its direction
is

\[
\boxed{
\text{extremely few inputs}
+
\text{extremely strong constraints}
\Longrightarrow
\text{dense geometry, algebra, and symmetry}
}
\]

In particular, the essential discrete parameter externally scannable in
this paper is \(N\), and

\[
M=\frac{N(N-1)}2
\]

is determined by \(N\). Names such as space, time, curvature, and
internal quantity are not additional parameters; they are observational
namings of the symmetry sectors that survive the strong constraints.

Accordingly, ``selecting a symmetry'' does not mean choosing an
arbitrary group from outside. It means that among the candidate
decompositions, the sectors that simultaneously satisfy the axiom system
and the self-consistency condition and close stably are selected, and
the observer reads their stabilizer / automorphism as a symmetry.

\[
\boxed{
\text{anonymous relational system}
\longrightarrow
\text{selection of allowed sectors by strong constraints}
\longrightarrow
\text{stable geometry and symmetry}
\longrightarrow
\text{readout as physical quantities}
}
\]

This distinction bears directly on the predictive power of the paper.
Even though multiple readout candidates exist, candidates that fail the
constraints are rejected, so this is not a construction that ``can
explain anything''. Rather, the fact that many mutually independent
structures must be satisfied simultaneously from the same few axioms
strengthens the falsifiability of this construction.

\section{2. A1: Complex Zero Closure}\label{a1-complex-zero-closure}

Let

\[
x=a+ib,
\qquad a,b\in\mathbb R^M
\]

Then

\[
0=x^Tx
=
a^Ta-b^Tb+2i\,a^Tb.
\]

Therefore

\[
\boxed{
a^Ta=b^Tb
}
\]

and

\[
\boxed{
a^Tb=0
}
\]

hold simultaneously.

\subsection{Theorem 1: Isometric Orthogonal
Two-Planes}\label{theorem-1-isometric-orthogonal-two-planes}

A nonzero solution of A1 requires, within the real relational space,

\[
\boxed{
\|a\|=\|b\|,
\qquad
a\perp b
}
\]

That is, A1 is not a mere ``zero sum'': it has isometric orthogonal
two-planes built in.

The global phase

\[
x\mapsto e^{i\theta}x
\]

is represented as a rotation within this plane.

\begin{center}\rule{0.5\linewidth}{0.5pt}\end{center}

\section{3. Observable and Unobservable
Axes}\label{observable-and-unobservable-axes}

Split the components of A1 into components \(r_a\) read as real
directions on the observation side and components \(s_\alpha\) carrying
\(i\) on the unobservable side.

\[
x=
(r_1,\ldots,r_p,
is_1,\ldots,is_q).
\]

Then A1 becomes

\[
\sum_{a=1}^{p}r_a^2
-
\sum_{\alpha=1}^{q}s_\alpha^2
=0.
\]

Hence, from the complex zero closure, an indefinite quadratic form of
signature \((p,q)\) appears automatically in the real display. The
signature here is not a metric signature imposed from outside, but the
count of \textbf{the number of axes read as real, \(p\), and the number
of axes read as unobservable imaginary, \(q\)}.

\[
\boxed{
Q_{p,q}
=
r^2-s^2
}
\]

In this sense the negative sign is not the result of imposing a
Minkowski metric from outside; it arises from the real-form display of
complex axes,

\[
\boxed{
(i s_\alpha)^2=-s_\alpha^2
}
\]

\begin{center}\rule{0.5\linewidth}{0.5pt}\end{center}

\section{\texorpdfstring{4. The \(R\) Axis Is Not a New Constant but a
Complex
Component}{4. The R Axis Is Not a New Constant but a Complex Component}}\label{the-r-axis-is-not-a-new-constant-but-a-complex-component}

As the minimal example, consider

\[
x=(x,y,z,iR)
\]

A1 is

\[
x^2+y^2+z^2-R^2=0.
\]

Therefore

\[
\boxed{
x^2+y^2+z^2=R^2
}.
\]

Here \(R\) is not a curvature radius added to the right-hand side.

It is a component inside the first axiom,

\[
\boxed{
x_R=iR
}
\]

If the observer does not read the \(R\) direction directly and reads
only \((x,y,z)\), the visible three directions satisfy the spherical
condition of radius \(R\).

Therefore \textbf{the coordinate value of a hidden complex axis is read
in visible space as a curvature radius}.

This is the central geometric interpretation of the paper.

\begin{center}\rule{0.5\linewidth}{0.5pt}\end{center}

\section{\texorpdfstring{5. Local Curvature Is a Readout of the \(R\)
Axis}{5. Local Curvature Is a Readout of the R Axis}}\label{local-curvature-is-a-readout-of-the-r-axis}

If a local section of visible space satisfies

\[
x^2+y^2+z^2=R^2
\]

the visible quadric is a sphere and the sectional curvature is

\[
\boxed{
K=\frac{1}{R^2}
}
\]

In general, on a \(d\)-sphere \(S^d(R)\),

\[
R_{abcd}
=
\frac1{R^2}
(g_{ac}g_{bd}-g_{ad}g_{bc}),
\]

\[
R_{ab}
=
\frac{d-1}{R^2}g_{ab},
\]

\[
\boxed{
\mathcal R
=
\frac{d(d-1)}{R^2}
}.
\]

Therefore, if \(R\) differs from place to place on the visible side,

\[
R=R(q)
\]

is \textbf{a local coordinate of the unobservable axis} and can
simultaneously be read on the visible side as a local curvature radius.

What matters is that

\[
\boxed{
R(q)\text{ has not been added as an external gravitational field}
}
\]

\(R\) is an internal component of A1.

\begin{center}\rule{0.5\linewidth}{0.5pt}\end{center}

\section{6. The Time Axis Works the Same
Way}\label{the-time-axis-works-the-same-way}

For

\[
x=(x,y,z,it)
\]

A1 is

\[
x^2+y^2+z^2-t^2=0.
\]

This is a Lorentz-type null cone.

Therefore the negative sign of time also appears as the real-form
display of an unobservable complex axis,

\[
\boxed{it}
\]

There is no need to assume a separate Minkowski metric here.

More generally, for

\[
x=(x,y,z,it,iR,iQ)
\]

we get

\[
\boxed{
x^2+y^2+z^2-t^2-R^2-Q^2=0
}.
\]

Therefore

\begin{itemize}
\tightlist
\item
  visible spatial axes,
\item
  the time axis,
\item
  the curvature axis \(R\),
\item
  the internal axis \(Q\),
\end{itemize}

need not be posited as separate fundamental principles. In particular,
the negative signs of \(t\), \(R\), and \(Q\) all come from the same
complex squaring rule

\[
\boxed{(iu)^2=-u^2}
\]

All are real-form displays of different complex directions of the same

\[
\boxed{
\sum_nx_n^2=0
}
\]

\subsubsection{Theorem 2: Lorentz Signature Arises from the Square of
Imaginary
Axes}\label{theorem-2-lorentz-signature-arises-from-the-square-of-imaginary-axes}

Let the observable axes be \(r_1,\ldots,r_p\) and the unobservable axes
\(u_1,\ldots,u_q\), and write the complex state as

\[
x=(r_1,\ldots,r_p,iu_1,\ldots,iu_q)
\]

A1 is

\[
\sum_{a=1}^{p}r_a^2+\sum_{\alpha=1}^{q}(iu_\alpha)^2=0
\]

that is,

\[
\boxed{
\sum_{a=1}^{p}r_a^2-\sum_{\alpha=1}^{q}u_\alpha^2=0
}
\]

Therefore the signature \((p,q)\) appearing in the real display is not
an independently introduced metric signature but is determined by
\textbf{the numbers of visible real axes and unobservable imaginary
axes}.

In particular, for \(p=3,q=1\),

\[
\boxed{
x^2+y^2+z^2-t^2=0
}
\]

and the Lorentz-type signature is a direct consequence of

\[
\boxed{(it)^2=-t^2}
\]

\(\square\)

In this sense the sign origins of the \(R\) axis and the \(t\) axis are
exactly identical.

\begin{center}\rule{0.5\linewidth}{0.5pt}\end{center}

\section{7. Conservation Groups Emerging from
A1}\label{conservation-groups-emerging-from-a1}

The complex linear conservation group of the quadratic form

\[
Q(x)=x^Tx
\]

is

\[
\boxed{
O(M,\mathbb C)
}.
\]

When the signature is chosen as \((p,q)\) in the real display, the real
conservation group is

\[
\boxed{
O(p,q)
}.
\]

The decomposition appearing in this six-component readout of the
construction is not \(3+1\) but

\[
\boxed{
3+3=(x,y,z)+(t,R,Q)
}
\]

That is, three axes on the visible side and three on the unobservable
side.

For the \((x,y,z,t)\) partial readout, the automorphism group of its
real quadratic form appears:

\[
\boxed{
O(3,1)
}
\]

The connected, orthochronous component is

\[
SO^+(3,1),
\]

and its double cover is

\[
\boxed{
Spin^+(3,1)\cong SL(2,\mathbb C)
}.
\]

\begin{center}\rule{0.5\linewidth}{0.5pt}\end{center}

\section{8. Null Cone, Curvature, and Spacetime Are Different Readouts
of the Same
A1}\label{null-cone-curvature-and-spacetime-are-different-readouts-of-the-same-a1}

A1 itself is always

\[
\sum_nx_n^2=0.
\]

But when the observer cannot read some axes directly, it appears as

\[
\sum_{\rm visible}r_a^2
=
\sum_{\rm hidden}s_\alpha^2
\]

If the right-hand side has just one axis, it can be read as a curvature
radius \(R\).

If the right-hand side includes the time axis, it can be read as a
Lorentz-type null structure.

Including several unobservable axes gives

\[
r^2=t^2+R^2+Q^2+\cdots
\]

Therefore

\[
\boxed{
\text{zero closure}
\to
\text{signature}
\to
\text{curvature}
\to
\text{spacetime}
}
\]

is not a sequence of different axioms but differences of decomposition,
projection, and readout of the same A1.

\begin{center}\rule{0.5\linewidth}{0.5pt}\end{center}

\section{9. Normalization Erodes This
Geometry}\label{normalization-erodes-this-geometry}

A1 is a homogeneous condition:

\[
Q(\lambda x)=\lambda^2Q(x).
\]

Therefore A1 does not kill the scale direction.

On the other hand, imposing

\[
x^\dagger x=1
\]

as a fundamental axiom fixes the radius of the Hermitian norm.

If the magnitudes of the unobservable complex axes are read on the
visible side as

\begin{itemize}
\tightlist
\item
  curvature radius,
\item
  clock scale,
\item
  internal scale,
\end{itemize}

then global normalization risks freezing those degrees of freedom.

Here one must not confuse

\[
x^Tx
\]

with

\[
x^\dagger x
\]

\[
\boxed{
x^Tx=0
}
\]

is the geometric zero closure.

\[
\boxed{
x^\dagger x=1
}
\]

is a section choice of the Hermitian norm.

\begin{center}\rule{0.5\linewidth}{0.5pt}\end{center}

\section{10. Complex Structure and Symplectic
Structure}\label{complex-structure-and-symplectic-structure}

In the real display

\[
x=(a,b)
\]

the complex structure is

\[
J=
\begin{pmatrix}
0&-I\\
I&0
\end{pmatrix},
\qquad
J^2=-I.
\]

For a compatible metric \(g\),

\[
\omega(u,v)=g(Ju,v)
\]

is an antisymmetric two-form.

Therefore

\[
\boxed{
\text{complex}
\to
\text{orthogonal}
\to
\text{symplectic}
}
\]

are not mutually independent extra assumptions; under the compatibility
condition they are different aspects of the same complex geometry.

\begin{center}\rule{0.5\linewidth}{0.5pt}\end{center}

\section{\texorpdfstring{11. A2: \(U^N=I\) Is Self-Closure of the
Action}{11. A2: U\^{}N=I Is Self-Closure of the Action}}\label{a2-uni-is-self-closure-of-the-action}

If

\[
\boxed{
U^N=I
}
\]

then every eigenvalue \(\lambda\) satisfies

\[
\lambda^N=1
\]

Therefore

\[
\boxed{
\lambda_m=e^{2\pi i m/N}
}.
\]

If the operator has exact order \(N\), then

\[
\boxed{
\langle U\rangle\cong\mathbb Z_N
}.
\]

Furthermore, the cyclotomic field

\[
\mathbb Q(\zeta_N)
\]

and

\[
\operatorname{Gal}
(\mathbb Q(\zeta_N)/\mathbb Q)
\cong
(\mathbb Z/N\mathbb Z)^\times
\]

are naturally attached.

A2 therefore provides simultaneously

\begin{itemize}
\tightlist
\item
  finite recurrence,
\item
  phase discretization,
\item
  cyclic groups,
\item
  cyclotomic structure.
\end{itemize}

\begin{center}\rule{0.5\linewidth}{0.5pt}\end{center}

\section{\texorpdfstring{12. Born-Type \(\cos^2/\sin^2\) Is a Readout of
Finite
Recurrence}{12. Born-Type \textbackslash cos\^{}2/\textbackslash sin\^{}2 Is a Readout of Finite Recurrence}}\label{born-type-cos2sin2-is-a-readout-of-finite-recurrence}

Imposing

\[
U^N=I
\]

on the two-channel rotation

\[
U(\theta)
=
\begin{pmatrix}
\cos\theta&-\sin\theta\\
\sin\theta&\cos\theta
\end{pmatrix}
\]

gives

\[
\theta=\frac{2\pi m}{N}.
\]

The squared projections onto the axes are

\[
\boxed{
\cos^2\theta,\qquad\sin^2\theta
}.
\]

Therefore Born-type squared weights can be generated, at least in
two-channel systems, as

\[
\boxed{
\text{finite-order phase recurrence}
\to
\text{squared readout}
}
\]

Born-type squared weights are not a starting axiom.

\begin{center}\rule{0.5\linewidth}{0.5pt}\end{center}

\section{\texorpdfstring{13. A3: The Same \(N\) Becomes the Vertex
Count}{13. A3: The Same N Becomes the Vertex Count}}\label{a3-the-same-n-becomes-the-vertex-count}

Let the \(N\) of A2 simultaneously be the number of geometric vertices.

Then the total number of pairwise relations is

\[
\boxed{
M=\frac{N(N-1)}2
}.
\]

Take each relation to be a scalar distance \(d_{ij}\) and assume all
distances are consistent as a simplex.

This establishes

\[
\boxed{
\text{number of phase divisions}
=
\text{number of vertices}
}
\]

A2 and A3 carry no independent integers; they are constrained to the
same \(N\).

This is a strong condition directly linking quantization and geometry.

\begin{center}\rule{0.5\linewidth}{0.5pt}\end{center}

\section{14. Simplex Closure Creates
Space}\label{simplex-closure-creates-space}

The pairwise distances of \(N\) vertices are the edge data of

\[
K_N
\]

If the distances are consistent, the embedding dimension and shape are
determined by the Gram matrix or the Cayley--Menger determinant.

Therefore

\[
\boxed{
\text{relational set}
\to
\text{distance geometry}
}
\]

holds.

Absolute positions are unnecessary.

Global

\begin{itemize}
\tightlist
\item
  translations,
\item
  rotations,
\item
  reflections,
\end{itemize}

do not change the distance data.

Hence they become redundancies in the relational description.

\begin{center}\rule{0.5\linewidth}{0.5pt}\end{center}

\section{15. The Simplex Also Generates a Chain
Complex}\label{the-simplex-also-generates-a-chain-complex}

The boundary operator of oriented simplices satisfies

\[
\boxed{
\partial^2=0
}
\]

Dualizing,

\[
\boxed{
d^2=0
}.
\]

Therefore the simplex is not merely a ``solid'': it generates

\begin{itemize}
\tightlist
\item
  homology,
\item
  cohomology,
\item
  discrete differential forms,
\item
  closed / exact structure.
\end{itemize}

This is the natural entry to the formal structure required by gauge
geometry.

\begin{center}\rule{0.5\linewidth}{0.5pt}\end{center}

\section{16. A4: Self-Consistency Selects Which Geometry Is
Realized}\label{a4-self-consistency-selects-which-geometry-is-realized}

Require

\[
\boxed{
X=\mathcal F(X)
}
\]

If \(\mathcal F\) is equivariant under a group \(G\), the conservation
group of a fixed point \(X_*\) is

\[
\boxed{
G_{X_*}
=
\{g\in G\mid gX_*=X_*\}
}.
\]

Therefore a symmetry selection

\[
\boxed{
G\to H
}
\]

occurs naturally.

The eigenvalues of the linearization

\[
\delta X'
=
D\mathcal F(X_*)\delta X
\]

select stable, neutral, and amplified directions.

Hence the phenomenon in which only a few directions become macroscopic
can be described as the spectral selection of self-consistent fixed
points.

\begin{center}\rule{0.5\linewidth}{0.5pt}\end{center}

\section{17. A2 and A4 Are Different Levels of
Self-Closure}\label{a2-and-a4-are-different-levels-of-self-closure}

\[
U^N=I
\]

is the finite recurrence of the whole operator.

\[
X=\mathcal F(X)
\]

is the fixed-point closure of states and geometry.

Extending to periodic fixed points

\[
\mathcal F^N(X)=X
\]

the two become different realizations of the same ``finite
self-recurrence''.

Therefore, taking

\[
\boxed{
\text{self-closure}
}
\]

as the superordinate principle, we can classify as one family:

\begin{itemize}
\tightlist
\item
  zero closure of quadratic forms,
\item
  finite closure of operators,
\item
  simplex closure of distances,
\item
  fixed-point closure of generating maps.
\end{itemize}

\begin{center}\rule{0.5\linewidth}{0.5pt}\end{center}

\section{\texorpdfstring{18. If \(3+1\) Emerges, Lorentz Is Not an
Additional
Axiom}{18. If 3+1 Emerges, Lorentz Is Not an Additional Axiom}}\label{if-31-emerges-lorentz-is-not-an-additional-axiom}

Suppose the complex-axis decomposition of A1 and the geometric selection
of A3/A4 stably select

\[
\boxed{
3\text{ visible real axes}
+
1\text{ unobservable imaginary axis}
}
\]

Writing the unobservable axis as \(it\), its square is automatically
\(-t^2\).

Then A1, in its real display, gives

\[
x^2+y^2+z^2+(it)^2=0
\]

that is,

\[
\boxed{x^2+y^2+z^2-t^2=0}
\]

The conservation group of this indefinite quadratic form is

\[
O(3,1)
\]

Therefore the unresolved point is neither the negative sign of Lorentz
nor whether three visible directions arise. The negative sign is
automatic from the square of \(it\).

Moreover, three-dimensionality itself is not treated as a mere accident
of numerical experiments. From the complex zero closure of the first
axiom arises the equal-norm orthogonal structure of real and imaginary
parts, and three directions are established by the orthogonal two-planes
and their independent normal direction. In addition, from the
consistency of phase closure and simplex geometry, the visible section
is represented as a three-dimensional ellipsoidal structure with three
mutually distinguishable principal axes.

Accordingly, the geometric path

\[
\boxed{
\text{zero closure}
\Longrightarrow
\text{equal-norm orthogonal structure}
\Longrightarrow
\text{orthogonal two-planes + normal}
\Longrightarrow
\text{3-dimensional ellipsoidal structure}
}
\]

is classified as the already-derived part of three-dimensionality.

In parallel, the author's published paper ``Zero Closure Was
Four-Dimensional'' numerically confirms that, of at most \(N-1\)
nontrivial principal axes, the top three directions \(A,B,C\) amplify
selectively, the middle directions remain almost stationary, and some
lower directions migrate to imaginary directions. Therefore

\[
\boxed{
\text{phase and closure geometry}
\Longrightarrow
\text{3-dimensional ellipsoidal structure}
}
\]

and

\[
\boxed{
\text{evolution of a self-consistent coherent parent}
\Longrightarrow
\text{spectral concentration and selective amplification onto 3 principal axes}
}
\]

are distinct grounds, yet point to the same three-dimensional visible
structure.

Hence we do not classify ``the principle generating three directions''
as unresolved. The geometric establishment of three directions and their
dominance and saturation are placed in the domain already derived and
confirmed in the self-papers.

What remains is the \textbf{elevation of this coincidence to a general
theorem of universality}. That is, for any sufficiently large \(N\), a
wide range of initial conditions, seeds, and admissible self-consistent
maps, to determine analytically the necessary and sufficient conditions
under which

\[
\boxed{
\text{3-dimensional ellipsoidal structure}
\Longleftrightarrow
\text{selective dominance of 3 principal axes}
}
\]

holds. In other words, the general analytic mechanism by which a
self-consistent coherent parent, a fermionic seed with odd harmonics,
and selective, inflation-like amplification preserving the zero closure
choose three anisotropic principal axes out of many relational
directions has not yet been determined.

Therefore the future task is not ``to create three directions'' but

\[
\boxed{
\text{to generalize the already reproduced spontaneous generation of three directions and to identify its selection principle analytically}
}
\]

\begin{center}\rule{0.5\linewidth}{0.5pt}\end{center}

\section{\texorpdfstring{19. Because of the \(R\) Axis, Local Curvature
Is Already Contained in
A1}{19. Because of the R Axis, Local Curvature Is Already Contained in A1}}\label{because-of-the-r-axis-local-curvature-is-already-contained-in-a1}

Conventional organizations tend to treat ``the generation of local
curvature'' as a separate problem.

But if A1 contains the unobservable component

\[
iR
\]

then in the visible three directions

\[
x^2+y^2+z^2=R^2
\]

Therefore

\[
\boxed{
R^{-2}
}
\]

is precisely the curvature of the visible sphere.

If locally

\[
R=R(q)
\]

the curvature radius of the visible section observed at each point
varies.

Therefore ``the place where gravity lives'' is not the addition of an
external force to the state; it already exists inside A1 as

\[
\boxed{
\text{the local readout of an unobservable complex axis}
}
\]

General Riemannian dynamics remains a derivation task, but \textbf{local
curvature as a geometric quantity need not be added as a new degree of
freedom}.

\begin{center}\rule{0.5\linewidth}{0.5pt}\end{center}

\section{20. Degrees of Freedom of the Complex Zero Closure, Readout
Resolution, and Two Routes to the Standard Model Gauge
Group}\label{degrees-of-freedom-of-the-complex-zero-closure-readout-resolution-and-two-routes-to-the-standard-model-gauge-group}

\begin{quote}
\textbf{Note on layers and selection.} The \(\mathbb C^6\), \(3+3\), and
\(3+2\) below do not denote a unique ontological number of axes fixed in
the fundamental relational system. The essential discrete parameter
given externally on the fundamental side is \(N\), and the number of
relational components is fixed as \(M=N(N-1)/2\). \(\mathbb C^6\) is one
effective representation reading the zero closure contracted into six
registers, and the complex-dimension count \(6-1=5\) after adopting that
representation is rigorous. Meanwhile, refinements like
\(Q^2=Q_1^2+Q_2^2+Q_3^2\) and re-contractions like \(r^2=t^2+R^2+Q^2\)
are other symmetric readouts of the same closure. The problem is
therefore not ``to derive the true number of axes uniquely from the
axioms'' but to classify the readout sectors, and their symmetries, that
close stably while satisfying the strong constraints simultaneously.
\end{quote}

The first axiom is read from the start as a single complex zero closure
on complex space

\[
\boxed{\sum_{n=1}^{6}X_n^2=0,\qquad X_n\in\mathbb C}
\]

Accordingly, the degree-of-freedom count here is not an argument
subtracting one real constraint from six real axes. Correctly,

\[
\boxed{\dim_{\mathbb C}\mathbb C^6-1=5}
\]

In the real display, the single complex constraint corresponds to two
real conditions, leaving 10 real dimensions out of 12 --- again five
complex dimensions.

Reading, as one observation map,

\[
\boxed{x^2+y^2+z^2=t^2+R^2+Q^2}
\]

the components \((x,y,z,it,iR,iQ)\) are ontologically all the same kind
of complex axis, and the distinctions visible/unobservable and
time/curvature/internal are attached on the readout side. In this
display, \(t\) can be placed as a dependent readout that is not read
directly,

\[
\boxed{t^2=x^2+y^2+z^2-R^2-Q^2}
\]

In the sector where this readout adopts \(t\) as a real clock quantity,

\[
\boxed{
x^2+y^2+z^2\ge R^2+Q^2
}
\]

is required, and on the equality surface

\[
\boxed{
x^2+y^2+z^2=R^2+Q^2
}
\]

we have \(t=0\). This is a boundary arising automatically from the
readout in question. At this stage the paper does not identify it with
the event horizon of general relativity; it is classified as a
\textbf{horizon-type boundary candidate}.

Therefore \(3+2\) is not a number introduced from outside to match the
Standard Model; it is one decomposition that reads, via an observation
map, the five complex degrees of freedom obtained by imposing one
complex zero closure on six complex axes.

\subsection{20.1 Coarse-Grained and Refined
Readouts}\label{coarse-grained-and-refined-readouts}

What matters is not to identify \(Q\) with the number of fundamental
components itself. \(Q\) is a register that reads the unobservable-side
closure quantity in contracted form; at a finer resolution it can be
expanded as

\[
\boxed{Q^2=Q_1^2+Q_2^2+Q_3^2}
\]

Therefore

\[
x^2+y^2+z^2-t^2=R^2+Q^2
\]

and

\[
\boxed{x^2+y^2+z^2-t^2=R^2+Q_1^2+Q_2^2+Q_3^2}
\]

are not two competing ontologies; they can be treated as a
coarse-graining/refinement relation reading the same zero closure at
different internal resolutions.

With this distinction, one must not confuse the number of relational
components counted by A3's

\[
M=\frac{N(N-1)}2
\]

with the number of readout axes \((r,t,R,Q_i,\ldots)\). Expanding
readout registers is not, by itself, an operation that changes \(M\).

\subsection{20.2 The Coarse-Grained Route to the Standard Model Gauge
Group}\label{the-coarse-grained-route-to-the-standard-model-gauge-group}

If the five complex degrees of freedom preserve a self-consistent
\(3\oplus2\) decomposition

\[
V=V_3\oplus V_2,
\qquad
\dim_{\mathbb C}V_3=3,
\qquad
\dim_{\mathbb C}V_2=2
\]

and each subspace is granted a Hermitian structure, the conservation
group is

\[
U(3)\times U(2)
\]

If, further, the conservation condition removing the physically
redundant global phase is realized as removal of the total determinant
phase, we obtain

\[
\boxed{
S(U(3)\times U(2))
\cong
\frac{SU(3)\times SU(2)\times U(1)}{\mathbb Z_6}
}
\]

Note that it is not necessary to identify \(V_3\) directly with visible
space \((x,y,z)\) and let color \(SU(3)\) act there. The central
principle of this paper is the separation of the ``complex zero closure
of the existence layer'' from the ``observational readout'', and the
target on which internal symmetries act must be identified separately as
a readout register.

Also, since in A1 the global phase

\[
X\mapsto e^{i\theta}X
\]

is treated as an unobservable common phase, the \(\det=1\) condition is
not a new independent axiom: it can very likely be derived as the
operation removing this global-phase redundancy from the physical
automorphism group. Rigorous treatment including the identification of
the discrete center remains, but the residual task is not ``assuming
det=1 from outside''; it is ``determining how far the global-phase
quotient of A1 uniquely fixes the global form of
\(S(U(3)\times U(2))\)''.

\subsection{20.3 The Refined Route Expanding the Internal
Triplet}\label{the-refined-route-expanding-the-internal-triplet}

In the refined readout

\[
R^2+Q_1^2+Q_2^2+Q_3^2
\]

the internal triplet \((Q_1,Q_2,Q_3)\) appears independently of the
three visible spatial axes. Preserving the complex Hermitian structure
of this triplet gives an internal \(U(3)\), and separating the global
phase gives an entry to \(SU(3)\).

In the readout collecting \((R,Q_1,Q_2,Q_3)\) as four complex
components, known group theory provides the alternative route

\[
U(4)\supset SU(4),
\qquad
SU(4)\supset SU(3)\times U(1)
\]

Meanwhile, the four components on the \((x,y,z,t)\) side admit a Lorentz
readout and a Euclidean-type readout; the double cover of the latter
exhibits

\[
Spin(4)\cong SU(2)\times SU(2)
\]

Hence the refined readout also has a known mathematical connection to
the Pati--Salam type

\[
SU(4)\times SU(2)\times SU(2)
\]

However, at the present stage this paper does not classify this as the
physical gauge group uniquely selected by A1--A4. What matters is that
the coarse-grained \(3+2\) route and the internal-triplet expansion
route arise from different readout resolutions of the same complex zero
closure.

\subsection{\texorpdfstring{20.4 The \(M=6\) Complex Zero Quadratic Form
and Known
Geometry}{20.4 The M=6 Complex Zero Quadratic Form and Known Geometry}}\label{the-m6-complex-zero-quadratic-form-and-known-geometry}

The projectivized zero set of the \(M=6\) A1 can be treated as the
nondegenerate complex quadric hypersurface \(Q^4\) in the complex
projective space \(\mathbb{CP}^5\). Quadrics of this kind connect with
Grassmann/Plücker geometry and with the accidental isomorphisms of
six-dimensional orthogonal groups. Therefore, by choosing different real
forms and observational readouts of the same complex zero closure, there
exist external mathematical routes connecting to \(SO(3,3)\), to the
\(SO(4,2)\) appearing in conformal geometry, and to
\(Spin(6)\cong SU(4)\).

This is not a claim to have proved the Standard Model directly from A1.
But it shows that the central proposition of this paper,

\[
\boxed{\text{one complex zero geometry}\longrightarrow\text{different readouts}\longrightarrow\text{different physical symmetries}}
\]

connects naturally to the known mathematics of complex quadric
hypersurfaces and their real forms. It is important as an external
mathematical cross-check route independent of the coarse-grained \(3+2\)
route.

\subsection{20.5 Derivation Stages at
Present}\label{derivation-stages-at-present}

\textbf{Directly obtained from A1}

\[
\boxed{\mathbb C^6\cap\{\sum X_n^2=0\}\Longrightarrow\dim_{\mathbb C}=5}
\]

and, as its observational readout,

\[
\boxed{x^2+y^2+z^2=t^2+R^2+Q^2}.
\]

\textbf{Already constructed and numerically confirmed, including the
self-papers}

The dominance of the three visible principal axes, the readout
decomposition including \(R/Q\), the refinement of internal structure,
and the particle-type classification by harmonic structure.

\textbf{Group-theoretically rigorous conditional part}

\[
3\oplus2+\text{preservation of the Hermitian decomposition}+\text{removal of the global phase}
\Longrightarrow
S(U(3)\times U(2))
\]

and

\[
S(U(3)\times U(2))
\cong
\frac{SU(3)\times SU(2)\times U(1)}{\mathbb Z_6}.
\]

\textbf{The remaining generalization and dynamics tasks} are to
generalize which \(3+2\) readout or internal triplet the rigorously
obtained five complex degrees of freedom select self-consistently ---
not to simply rebuild readouts already constructed. They are to
determine which readout resolution and which stabilizer are universally
selected by self-consistent dynamics, how local gauge connections arise,
and whether chirality, hypercharge, and anomaly cancellation close
within the same derivation chain.

\section{21. Revised Derivation Map}\label{revised-derivation-map}

The top-level principle for reading the derivation map is

\[
\boxed{
N
\Longrightarrow
M=\frac{N(N-1)}2
\Longrightarrow
\text{anonymous relational system}
\Longrightarrow
\text{sector selection by strong constraints}
\Longrightarrow
\text{readout of symmetry and geometry}
}
\]

Accordingly, \(3+3\), \(3+2\), the internal \(Q\) triplet, etc. are not
mutually competing ``fundamental axis counts''; they are classified as
readouts of different resolutions and symmetry sectors of the same
strongly constrained relational system. The admissibility of multiple
readouts is not, in itself, the addition of free parameters.

{\def\LTcaptype{none} % do not increment counter
\begin{longtable}[]{@{}
  >{\raggedright\arraybackslash}p{(\linewidth - 4\tabcolsep) * \real{0.3333}}
  >{\raggedright\arraybackslash}p{(\linewidth - 4\tabcolsep) * \real{0.3333}}
  >{\raggedright\arraybackslash}p{(\linewidth - 4\tabcolsep) * \real{0.3333}}@{}}
\toprule\noalign{}
\begin{minipage}[b]{\linewidth}\raggedright
Structure
\end{minipage} & \begin{minipage}[b]{\linewidth}\raggedright
Origin
\end{minipage} & \begin{minipage}[b]{\linewidth}\raggedright
Status
\end{minipage} \\
\midrule\noalign{}
\endhead
\bottomrule\noalign{}
\endlastfoot
Anonymity / strong-constraint sector selection & \(N\to M=N(N-1)/2\);
intersection of zero closure, finite recurrence, simplex, harmonics,
self-consistency & Multiplicity of readouts treated as sector selection,
not arbitrariness \\
Complex zero quadratic form & A1 & Axiom \\
Isometric orthogonal two-planes & A1 & Rigorous \\
Complex phase rotation & A1 & Rigorous \\
\(O(M,\mathbb C)\) & Conservation group of A1 & Rigorous \\
Indefinite signature \(O(p,q)\) & Real display of unobservable imaginary
axes \(iu\) & Rigorous \\
Null cone & Real display of A1 & Rigorous \\
\(x^2+y^2+z^2=R^2\) & Including \(iR\) in A1 & Rigorous \\
Curvature radius \(R\) & Visible-sphere readout & Rigorous \\
\(K=1/R^2\) & Spherical geometry & Rigorous \\
Negative sign of \(t\) & Reading the unobservable axis as \(it\), so
\((it)^2=-t^2\) & Rigorous \\
Signature \((3,3)\) of the six-component readout & Real display of
\((x,y,z,it,iR,iQ)\) & Rigorous \\
\(O(3,3)\) & Conservation group of the \((3,3)\) quadratic form of the
six-component readout & Rigorous \\
Lorentz \(O(3,1)\) & Spacetime partial readout of \((x,y,z,t)\) &
Conditionally rigorous \\
Distinguishability of 3 visible principal axes & Anisotropic ellipsoid
\(a\neq b\neq c\) & Rigorous within the range where self-consistent
principal-axis solutions hold \\
Spin\((3,1)\) & Lorentz lift & Conditionally rigorous \\
Conformal structure & Null cone / scale & Strong known connection \\
Symplectic \(J\) & Real display of the complex structure & Rigorous
entry \\
\(\mathbb Z_N\) & A2 & Rigorous \\
Cyclotomic eigenvalues & A2 & Rigorous \\
Galois symmetry & Cyclotomic field & Rigorous \\
Born-type squared weights & A2 + two-channel projection & Rigorous \\
\(K_N\) & A3 & Rigorous \\
\(S_N\) & Anonymous vertices & Rigorous \\
Simplex distance geometry & A3 & Rigorous \\
\(\partial^2=0\) & Simplex chain complex & Rigorous \\
Cohomology & A3 & Rigorous \\
Stabilizer & A4 & Rigorous \\
\(G\to H\) & A4 & Conditionally rigorous \\
3-dimensional ellipsoidal structure & A1: equal-norm orthogonal
structure → orthogonal two-planes + normal & Already derived \\
Selective dominance of 3 principal axes & A4 / \texttt{make\_parent}
self-consistent spectrum & Numerically confirmed in the self-papers;
universality and necessary-and-sufficient conditions remain to be
generalized \\
Complex six-axis zero closure & A1 on \(\mathbb C^6\): one complex
constraint & \(\dim_{\mathbb C}=5\) derived rigorously \\
Coarse-grained readout & \(x^2+y^2+z^2=t^2+R^2+Q^2\) & Derived as an
observational display of the same complex zero closure \\
Refined internal readout & \(Q^2=Q_1^2+Q_2^2+Q_3^2\) & Consistent with
the readout conventions of the self-papers; nested with the
coarse-grained readout \\
Independent degrees of freedom & Six complex dof \(-\) one complex zero
closure & Rigorous as \(6-1=5\) complex dof; \(3+2\) is one
observational decomposition of it \\
\(U(3)\times U(2)\) & Preservation of the Hermitian \(3\oplus2\)
decomposition of the 5 complex dof & Conditionally rigorous; the
selection rule as a self-consistent stabilizer is the generalization
task \\
\(S(U(3)\times U(2))\) & The above + removal of the global-phase
redundancy of A1 & Conditionally rigorous; rigorous treatment of the
global form including the discrete center remains \\
Standard Model gauge Lie algebra & Lie algebra of \(S(U(3)\times U(2))\)
& \(\mathfrak{su}(3)\oplus\mathfrak{su}(2)\oplus\mathfrak{u}(1)\):
conditionally rigorous \\
SM global gauge group & \(S(U(3)\times U(2))\) &
\([SU(3)\times SU(2)\times U(1)]/\mathbb Z_6\): conditionally
rigorous \\
Internal \(Q\) triplet & Refined readout of \((Q_1,Q_2,Q_3)\) & Entry to
an internal \(SU(3)\) candidate independent of visible \((x,y,z)\) \\
\(SU(4)\) route & Four-complex-component readout of \((R,Q_1,Q_2,Q_3)\)
& Known connections to \(Spin(6)\cong SU(4)\) and
\(SU(4)\supset SU(3)\times U(1)\) \\
Pati--Salam-type route & Internal 4 components +
\(Spin(4)\cong SU(2)\times SU(2)\) & Second, externally mathematical
route; physical selection not yet generalized \\
Complex quadric \(Q^4\subset\mathbb{CP}^5\) & Projective zero set of the
\(M=6\) A1 & Strong connection to known mathematics; unified cross-check
route across real forms and readouts \\
Bose/Fermi/mixed sectors & Odd/even harmonic ratio, \(Z_2\) parity,
single/double cover & Derived, constructed, and numerically confirmed in
the self-papers; mass identification and general spin-statistics
correspondence remain \\
General Riemannian dynamics & \(R(q)\) + connection & Not yet derived \\
Chirality/hypercharge/anomaly & SM representations & Not yet derived \\
\end{longtable}
}

\begin{center}\rule{0.5\linewidth}{0.5pt}\end{center}

\subsubsection{21.1 The Double Derivation Branch Added in This
Revision}\label{the-double-derivation-branch-added-in-this-revision}

This revision retains the coarse-grained route

\[
\mathbb C^6\xrightarrow{\sum X_n^2=0}\dim_{\mathbb C}=5\xrightarrow{3\oplus2}S(U(3)\times U(2))
\]

while adding a second route that refines the internal readout to

\[
Q^2=Q_1^2+Q_2^2+Q_3^2
\]

This makes it unnecessary to let \(SU(3)\) act directly on the three
visible spatial axes; it can be read as a symmetry acting on the
internal triplet.

Furthermore, the \(M=6\) A1 has a known mathematical connection to the
complex quadric hypersurface \(Q^4\subset\mathbb{CP}^5\), providing an
independent route through \(Spin(6)\cong SU(4)\). The object of future
verification is therefore not ``constructing
\(SU(3)\times SU(2)\times U(1)\) from one accidental \(3+2\)
decomposition'', but \textbf{to which stabilizer the multiple readout
routes over the same complex zero geometry converge}.

\begin{center}\rule{0.5\linewidth}{0.5pt}\end{center}

\section{22. The Problems That Truly
Remain}\label{the-problems-that-truly-remain}

By reading the first axiom correctly, the following cease to be
``problems requiring new structure'':

\begin{itemize}
\tightlist
\item
  curvature radius,
\item
  the geometric location of local curvature,
\item
  the negative sign of the time axis,
\item
  the null cone,
\item
  indefinite quadratic forms,
\item
  the entry to the Lorentz group,
\item
  the scale degree of freedom.
\end{itemize}

In particular, \(t\), \(R\), and \(Q\) are unobservable complex axes of
the same kind, and their sign differences arise solely from

\[
\boxed{(iu)^2=-u^2}
\]

The remaining central problems narrow down as follows.

\subsection{Statistical Structure Taken as Derived
Self-Hypotheses}\label{statistical-structure-taken-as-derived-self-hypotheses}

The construction obtained in ``The Periodic Table of Waves v2'' is
adopted in this paper as the derivation chain

\[
\text{odd/even harmonics}
\longrightarrow
T_{\lambda_0/2}\psi=\pm\psi
\longrightarrow
Z_2\text{ parity}
\longrightarrow
\text{single/double cover}
\]

Furthermore, with the odd- and even-harmonic powers
\(P_{\rm odd},P_{\rm even}\),

\[
\boxed{
r=\frac{P_{\rm odd}}{P_{\rm odd}+P_{\rm even}}
}
\]

gives the exchange/reflection rate. This continuous quantity
distinguishes the pure-even endpoint, the pure-odd endpoint, and the
mixed sector in between.

Accordingly, in this paper,

\[
\boxed{
\begin{array}{ll}
P_{\rm odd}=0 &\rightarrow \text{Bose-type endpoint},\\[2mm]
P_{\rm even}=0 &\rightarrow \text{Fermi-type endpoint},\\[2mm]
P_{\rm odd}P_{\rm even}>0 &\rightarrow \text{mixed (Ermion-type) sector}
\end{array}
}
\]

is not a classification newly assumed from the external spin-statistics
theorem, but \textbf{a derived result, as a self-hypothesis, obtained
from the author's published model}.

Therefore the future task is not to ``generate'' the Bose/Fermi
structure. What remains is to compare and formalize, under more general
conditions, the correspondence between this harmonic-parity-derived
statistical structure and the spin-statistics structure of standard
quantum field theory.

\subsection{Problem 1}\label{problem-1}

For the spontaneous generation of the three anisotropic visible
principal axes already reproduced with \texttt{make\_parent}, identify
the generating principle analytically and generalize it beyond the
specific numerical construction. In particular, determine which of the
following are necessary and sufficient for three-direction selection: a
self-consistent coherent parent, a fermionic seed with odd harmonics,
zero closure, and selective amplification.

\subsection{Problem 2}\label{problem-2}

Among the multiple readout resolutions admitted by the complex zero
closure, generalize which Hermitian decomposition and stabilizer are
universally selected by self-consistent dynamics. This is not the
problem of deriving the \(3+2\) degrees of freedom anew; it is the
problem of determining the selection rule among the existing routes: the
coarse-grained \(3\oplus2\), the internal \(Q\) triplet, and the
\(SU(4)\)-type refinement.

\subsection{Problem 3}\label{problem-3}

Can general Riemann curvature and geodesic dynamics be derived from the
self-consistent variation of the unobservable axis

\[
R(q)
\]

\subsection{Problem 4}\label{problem-4}

Can local gauge connections / curvature be derived from simplex cochains
and phase data?

\subsection{Problem 5}\label{problem-5}

The Bose/Fermi/mixed classification itself is already derived,
constructed, and numerically confirmed in the self-papers. What remains
is to generalize the classification by odd/even harmonic ratio and
\(Z_2\) parity to arbitrary sufficiently large \(N\) and wide
interaction conditions, and to close the correspondence with the mass
identification and mass hierarchy of real particles and with the
spin-statistics structure of standard quantum field theory.

\subsection{Problem 6}\label{problem-6}

Can the chirality, hypercharge, and anomaly cancellation of the Standard
Model be reproduced from stabilizer representations?

\begin{center}\rule{0.5\linewidth}{0.5pt}\end{center}

\section{23. Falsifiability}\label{falsifiability}

The strong form of this construction is falsified by any of the
following.

\begin{enumerate}
\def\labelenumi{\arabic{enumi}.}
\tightlist
\item
  The self-consistent fixed points of A1--A4 fail to reproduce the
  observed three-principal-axis structure or the admissible readout
  sectors.
\item
  Identifying the \(N\) of A2 with the vertex count of A3 eliminates all
  solutions.
\item
  Local variation of the \(iR\) axis is inconsistent with the geodesic /
  curvature readout.
\item
  At large \(N\) the three-principal-axis selection disappears, reducing
  to a mere finite-size effect.
\item
  The automorphism group of the five-degree-of-freedom structure fails
  to reproduce the internal symmetries and representation content of the
  Standard Model.
\end{enumerate}

\begin{center}\rule{0.5\linewidth}{0.5pt}\end{center}

\section{24. Conclusion}\label{conclusion}

The central principle of this paper is, throughout,

\[
\boxed{
\sum_nx_n^2=0
}
\]

The curvature radius \(R\) is not a quantity added to the right-hand
side.

When the complex component

\[
\boxed{
x_R=iR
}
\]

is displayed in real form, it merely looks like

\[
\sum_{\rm visible}x_a^2=R^2
\]

Similarly, the time axis, as the complex direction

\[
it
\]

produces

\[
r^2-t^2=0
\]

Here \(t\), \(R\), and \(Q\) are essentially unobservable complex axes
of the same kind, and which axis is called time, curvature, or internal
quantity is determined by which observation map reads it and how. The
difference of names itself therefore need not be derived uniquely from
the axioms.

Hence

\[
\boxed{
\text{curvature}
,\quad
\text{time signature}
,\quad
\text{null structure}
}
\]

are not additions external to the first axiom.

They all emerge from

\[
\boxed{
\text{the observable/unobservable axis decomposition of the complex zero closure}
}
\]

Furthermore, when the six directions are read as

\[
(x,y,z,it,iR,iQ)
\]

the first axiom gives

\[
\boxed{
x^2+y^2+z^2=t^2+R^2+Q^2
}
\]

In this six-component readout the geometry is thus expressed not as
\(3+1\) but as \(3+3\), and the \(3+1\) of observed spacetime is a
partial readout of it. When the visible side becomes an anisotropic
ellipsoid, its three principal axes appear autonomously as mutually
distinguishable eigendirections.

The second axiom

\[
U^N=I
\]

provides the finite self-closure of the action;

the third condition, simplex closure, converts phase vertices into
distance geometry;

and the fourth condition, self-consistency, selects the fixed points and
symmetries realized among them.

The few principles therefore compress to

\[
\boxed{
\text{zero closure}
+
\text{finite recurrence}
+
\text{distance closure}
+
\text{self-consistency}
}
\]

Under this compression, many symmetries of theoretical physics can be
reclassified not as independent axioms but as

\[
\boxed{
\text{automorphisms of a self-closing complex relational geometry}
}
\]

The most important task is not to add further symmetry groups. It is

\[
\boxed{
\text{to classify completely the fixed points and stabilizers of A1--A4}
}
\]

\begin{center}\rule{0.5\linewidth}{0.5pt}\end{center}

\section{References}\label{references}

External literature is primary; only the three primary sources of
results derived and confirmed as self-hypotheses are self-cited,
minimally.

\begin{enumerate}
\def\labelenumi{\arabic{enumi}.}
\item
  E. W. Weisstein, ``Sphere,'' \emph{MathWorld}. Sphere geometry and
  curvature references.\\
  https://mathworld.wolfram.com/Sphere.html
\item
  E. W. Weisstein, ``Cayley-Menger Determinant,'' \emph{MathWorld}.\\
  https://mathworld.wolfram.com/Cayley-MengerDeterminant.html
\item
  M. Cvetič and L. Lin, ``The Global Gauge Group Structure of F-theory
  Compactification with U(1)s,'' arXiv:1706.08521 (2017).\\
  https://arxiv.org/abs/1706.08521
\item
  N. Raghuram, W. Taylor and A. P. Turner, ``General F-theory models
  with tuned \((SU(3)\times SU(2)\times U(1))/\mathbb Z_6\) symmetry,''
  arXiv:1912.10991 (2019).\\
  https://arxiv.org/abs/1912.10991
\item
  F. Klinker, ``An explicit description of \(SL(2,\mathbb C)\) in terms
  of \(SO^+(3,1)\) and vice versa,'' arXiv:1712.02168 (2017).\\
  https://arxiv.org/abs/1712.02168
\item
  J. T. Wheeler, ``Weyl geometry,'' arXiv:1801.03178 (2018).\\
  https://arxiv.org/abs/1801.03178
\item
  O. Macia and Y. Nagatomo, ``Holomorphic isometric embeddings of
  complex Grassmannians into quadrics: The general case,'' \emph{Kyoto
  Journal of Mathematics} 66(1), 67--85 (2026).\\
  https://doi.org/10.1215/21562261-2024-0033
\item
  Standard differential geometry result: a \(d\)-sphere of radius \(R\)
  has constant sectional curvature \(1/R^2\), Ricci curvature
  \((d-1)g/R^2\), and scalar curvature \(d(d-1)/R^2\).
\item
  Standard Lorentz geometry result: the real quadratic form of signature
  \((3,1)\) is preserved by \(O(3,1)\); its proper orthochronous spin
  double cover is \(SL(2,\mathbb C)\).
\item
  Standard algebraic topology result: the simplicial boundary operator
  satisfies \(\partial^2=0\).
\end{enumerate}

\begin{center}\rule{0.5\linewidth}{0.5pt}\end{center}

\subsection{Claim Strength}\label{claim-strength}

\subsubsection{Rigorous}\label{rigorous}

\begin{itemize}
\tightlist
\item
  \(x^Tx=0\Rightarrow \|a\|=\|b\|,\ a\perp b\)
\item
  That A1 containing the complex axis \(iR\) gives \(r^2=R^2\) on the
  visible side
\item
  That A1 containing the complex axis \(it\) gives the Lorentz-type
  negative sign on the visible side
\item
  \(K=1/R^2\) for visible spherical sections
\item
  \(\mathbb Z_N\) and cyclotomic eigenvalues from A2
\item
  \(K_N,S_N,\partial^2=0\) from A3
\item
  Stabilizers from A4 + equivariance
\end{itemize}

\subsubsection{Conditionally Rigorous}\label{conditionally-rigorous}

\begin{itemize}
\tightlist
\item
  Lorentz / Spin / conformal in the \((x,y,z,t)\) partial readout of
  \(3+3\)
\item
  \(S(U(3)\times U(2))\) and the Standard Model global gauge group after
  the \(3+2\) selection
\item
  The construction reading \(R(q)\) as a local visible curvature radius
\end{itemize}

\subsection{\texorpdfstring{The Only Externally Specified Parameter Is
the Discrete Integer
\(N\)}{The Only Externally Specified Parameter Is the Discrete Integer N}}\label{the-only-externally-specified-parameter-is-the-discrete-integer-n}

An important feature of this construction is that, apart from the parts
not yet derived or generalized, no explicit continuous tuning parameters
for fitting phenomena are introduced.

The only structural parameter specified from outside is

\[
\boxed{N\in\mathbb N}
\]

This \(N\) is not an ordinary coupling constant or fitting parameter. In
this construction the same integer \(N\) simultaneously binds

\[
\boxed{
U^N=I,\qquad
|V|=N,\qquad
M=\frac{N(N-1)}2
}
\]

specifying the degree of finite-order recurrence, the number of discrete
divisions of the phase circle, the number of vertices of the closure
geometry, and the total number of pairwise relations.

Therefore \(N\) is not a value tuned continuously to observations; it is
a \textbf{discrete structural degree} specifying the finiteness, closure
degree, and relational-space scale of the system itself.

Meanwhile, the quantities appearing in this paper and the three
self-referenced papers,

\[
R,\quad t,\quad Q,\quad
r,\quad
P_{\rm odd},\quad P_{\rm even},
\]

the principal-axis lengths, the clock field, the mass readout, the
reflectance, the covering degree, and the Bose/Fermi/mixed
classification, are not entered as independent free parameters; they are
treated as quantities arising from the zero closure, finite order,
simplex relations, harmonic structure, self-consistency, and observation
maps.

At the present stage, therefore, the model is characterized as a

\[
\boxed{
\text{one-discrete-parameter structural model}
}
\]

that is, \textbf{a structural model whose only externally specified
degree of freedom is the integer \(N\)}.

This is not merely a matter of ``having few parameters''. It means that
no individual masses, coupling constants, curvature scales, statistical
mixing rates, or anisotropies are added as separate tuning knobs to
obtain the symmetries, geometry, statistics, and readout structures the
construction sets out to explain.

Furthermore, if it is shown in the future that \(N\) itself is selected
uniquely or discretely by the self-consistency or closure conditions,
then

\[
\boxed{
\text{external free input}=0
}
\]

becomes possible. This, however, is not treated as derived at present.

\subsubsection{\texorpdfstring{The Remaining Generalization Task
Concerning
\(N\)}{The Remaining Generalization Task Concerning N}}\label{the-remaining-generalization-task-concerning-n}

In the present theory \(N\) is the only externally specified integer.
One of the remaining fundamental problems is therefore whether

\[
\boxed{
\text{why this universe / observation sector selects a particular }N
}
\]

can be derived from self-consistency, finite-order closure, or the
structure of the large-\(N\) limit.

This is not a task of adding a new continuous parameter to the derived
structure; it is \textbf{the task of finding the selection rule for the
single remaining discrete input \(N\)}.

\subsubsection{Derived, Constructed, and Numerically Confirmed in the
Self-Papers}\label{derived-constructed-and-numerically-confirmed-in-the-self-papers}

\begin{itemize}
\tightlist
\item
  \textbf{The 3-dimensional ellipsoidal structure and the dominance of
  three principal axes}\\
  Three-dimensionality is classified as already derived as a
  3-dimensional ellipsoidal structure, from the equal-norm orthogonal
  structure arising from the zero closure, the orthogonal two-planes
  with independent normal, and the consistency of phase closure and
  simplex geometry. Furthermore, ``Zero Closure Was Four-Dimensional''
  numerically confirms that of at most \(N-1\) nontrivial principal axes
  the top three directions \(A,B,C\) amplify selectively, the middle
  directions remain nearly stationary, and some lower directions migrate
  to imaginary directions. The establishment of three directions, the
  anisotropy, and the dominance and saturation in three directions are
  themselves not ``underived''.\\
  \textbf{Remaining task:} to elevate the coincidence of geometric
  three-dimensionality and dynamical three-axis dominance to a general
  theorem for arbitrary sufficiently large \(N\) and wide initial
  conditions.
\end{itemize}

The following are not items newly proved within this paper from A1--A4
alone, but are already derived, constructed, and numerically confirmed
as self-hypotheses in the author's three published papers. They are
therefore not classified as ``underived''.

\begin{itemize}
\item
  \textbf{Spontaneous generation of three directions}\\
  The construction in which, from a self-consistent coherent parent via
  \texttt{make\_parent}, a fermionic seed with odd harmonics, zero
  closure, and selective, inflation-like amplification, three
  anisotropic visible principal axes come to dominate among many
  relational directions, has been realized.\\
  \textbf{Remaining task:} the analytic elucidation of why three
  directions are selected, the extraction of necessary and sufficient
  conditions, and generalization to arbitrary sufficiently large \(N\).
\item
  \textbf{Bose/Fermi/mixed (Ermion) classification}\\
  The \(Z_2\) parity under half-period translation of odd and even
  harmonics, the reflectance determined by the odd-harmonic ratio \[
  r=\frac{P_{\rm odd}}{P_{\rm odd}+P_{\rm even}},
  \] and the Bose/Fermi/mixed classification through single/double
  covers and the antipodal two-point structure are derived, constructed,
  and numerically confirmed in ``The Periodic Table of Waves v2''.\\
  \textbf{Remaining task:} complete mass identification with real
  particle species, a closed derivation of the mass hierarchy, and the
  establishment of a general correspondence with the spin-statistics
  structure of standard quantum field theory.
\item
  \textbf{The two-layer structure of mass, clock field, and gravity
  readout}\\
  The construction separating the wave-existence layer from the
  observational readout layer, reading mass as the value of the clock
  field and gravity as its gradient, and its separation from gauge-type
  readout, are derived and numerically confirmed as self-hypotheses in
  ``Two-Layer Separation of Waves and Fields''.\\
  \textbf{Remaining task:} fixing the mass invariant, generalization to
  the real mass hierarchy, extension of the gravity readout to general
  variable-curvature systems, and full dynamicization.
\item
  \textbf{The readout-level distinction of \(t,R,Q\)}\\
  \(t,R,Q\) are fundamentally unobservable complex axes of the same
  kind; what is read as state updating is named time \(t\), what is read
  as the curvature scale of visible geometry is named \(R\), and what is
  read as an internal quantity is named \(Q\). This is an identification
  on observation maps, not a ``problem of deriving three kinds of axes
  from the axioms''.\\
  \textbf{Remaining task:} not a proof of principled role
  differentiation, but refining, as needed, the correspondence between
  each observation map and measured quantities.
\end{itemize}

\subsubsection{Not Yet Derived or Generalized, Even Including This Paper
and the
Self-Papers}\label{not-yet-derived-or-generalized-even-including-this-paper-and-the-self-papers}

Only what truly remains outside the closed derivation chain is placed
here.

\begin{itemize}
\item
  \textbf{General classification of readout sectors from the fundamental
  \(M\)-component system}

  Apart from the three-direction generation by the existing
  \texttt{make\_parent} construction, to obtain the selection of a
  specific number of axes as a general theorem from A1--A4 alone.
\item
  \textbf{Complete dynamics of general variable curvature}\\
  Not the local curvature structure itself, which reads \(R\) as an
  unobservable direction of the complex square closure, but closed,
  complete dynamics for arbitrarily varying curvature fields.
\item
  \textbf{Complete derivation of local gauge dynamics}\\
  Beyond the gauge-type readouts and selection rules already obtained,
  to close local gauge connections and their complete dynamics generally
  from the axiom system.
\item
  \textbf{Complete derivation of chirality / hypercharge / anomaly
  cancellation}\\
  To close uniquely the correspondence candidates with the derived
  winding-number, parity, and internal-quantity structures, up to the
  chirality, hypercharge, and anomaly cancellation conditions of the
  Standard Model.
\item
  \textbf{The selection rule for \(N\)}\\
  To derive whether the single remaining externally specified structural
  parameter \(N\) is selected autonomously from self-consistency,
  finite-order closure, the large-\(N\) limit, or similar. If this
  closes, the model has zero external free input.
\end{itemize}

\subsection{Self-References}\label{self-references}

\begin{itemize}
\tightlist
\item
  Noriaki Kihara, ``The Periodic Table of Waves v2 --- Particle
  Classification by Winding-Number Address and Observation Clock, and
  the Unification of Mass, Lifetime, and Splitting via the Clock Field
  \(\omega(x)\)'', Zenodo, Concept DOI: 10.5281/zenodo.21830706.
  \textbf{Used in this paper as the source of previously derived results
  as self-hypotheses.}
\item
  Noriaki Kihara, ``Two-Layer Separation of Waves and Fields ---
  Unification of Gauge Fields and Gravitational Fields via a Universal
  Field-Readout Function'', Zenodo, DOI: 10.5281/zenodo.21832257.
  \textbf{Used in this paper as the source of previously derived results
  as self-hypotheses.}
\item
  Noriaki Kihara, ``Zero Closure Was Four-Dimensional --- `Central
  Projection' Survives Even in the Complex World'', Zenodo, Version DOI:
  10.5281/zenodo.21902806, Concept DOI: 10.5281/zenodo.21902805.
  \textbf{Used in this paper as the source of the previously derived
  results, as self-hypotheses, on three-direction spectral
  concentration, anisotropic principal-axis generation, and
  inflation-like expansion.}
\end{itemize}

\end{document}
